\documentclass[profile=standard,twoside=false]{ipara-handbook}

\usetikzlibrary{arrows.meta,calc,fit,positioning,backgrounds}
\pgfdeclarelayer{background}
\pgfdeclarelayer{foreground}
\pgfsetlayers{background,main,foreground}

\handbooksetup
  {Attention 与 Transformer：信息路由、架构组合与高效实现}
  {人工智能 · Area 60 · Leaf 06}
  {Attention · Transformer · FlashAttention · Triton}

\begin{document}

\begin{titlepage}
  \centering
  \vspace*{0.15\textheight}
  {\Huge\bfseries\sffamily Attention 与 Transformer\par}
  \vspace{0.8em}
  {\Large 从内容路由到 FlashAttention：同一数学语义的高效实现\par}
  \vspace{1.8em}
  {\small 技术快照：2026-08-22\par}
  \vspace{1.2em}
  {\small\href{../../40_复试/10_人工智能与机器学习/60_深度学习/06_Attention与Transformer/README.md}{返回 Topic Landing Page}\par}
\end{titlepage}

\tableofcontents

\section{本册定位：Attention 是信息路由机制，Transformer 是架构组合}

序列建模的核心困难之一是：一个位置当前需要的信息可能来自很远的位置，而且“该看谁”并不是固定的。卷积用固定局部窗口汇总邻域，循环网络把历史压入递归状态；Attention 采用另一条路线：\textbf{让当前位置根据内容产生查询，并对候选信息动态分配权重。}

\begin{mentalmodel}{Mother Question｜怎样按内容动态路由信息}
给定一组表示，怎样让每个 query 位置根据自己的当前内容，对一组 key 计算相关性，再用这些相关性汇总对应 value；随后又怎样把这种路由机制与残差、归一化、前馈网络和位置信息组合成可堆叠的 Transformer？
\end{mentalmodel}

这先确定 Attention/Transformer 的数学与架构语义。只有在这套语义固定以后，才进入第二条由实现瓶颈触发的分支：当序列很长时，怎样在保持标准 Attention 数学语义的同时，避免大型中间矩阵反复写入和读取 device global memory（高端加速器上通常由 HBM 承载）？FlashAttention 回答的是这个实现问题，而不是重新定义 Attention。

本册有三层责任：

\begin{enumerate}
  \item \textbf{语义层}：Attention 为什么需要 $Q,K,V$，权重怎样产生，self/cross/causal/multi-head 分别改变什么；
  \item \textbf{架构层}：Attention 怎样进入 Transformer block，以及 positional information 为什么不可缺；
  \item \textbf{实现层}：朴素 Attention 的数据流为什么昂贵，online softmax 怎样使 tiled exact attention 成为可能，FlashAttention 与 Triton 又如何把数学结构映射到 GPU。
\end{enumerate}

通用 reverse-mode AD、JVP/VJP 不在本册重新定义；本册只调用 Area 50 AutoDiff 给出的 VJP 接口，并负责 Attention 自己的 backward 组合。

\begin{boundarytable}
\boundarytablehead
Attention & $\neq$ & Transformer
  & Attention 是动态信息路由算子；Transformer 是围绕 attention、FFN、residual、normalization、position 等构造的架构族
  & 会把一个算子误当成完整网络架构\\
Transformer & $\neq$ & LLM
  & Transformer 是 architecture family；LLM 还包含训练目标、数据、优化、tokenization、decoding 与系统实现
  & 会把“使用 Transformer”误当成“已经定义语言模型生命周期”\\
FlashAttention & $\neq$ & Approximate Attention
  & FlashAttention 重新安排精确 Attention 的计算与数据搬运；它不靠低秩/稀疏近似改变标准输出
  & 会把性能优化误解成数学近似\\
FlashAttention & $\neq$ & KV Cache
  & FlashAttention 优化一次 attention kernel 的计算/内存流；KV cache 复用自回归推理中历史 token 的 K/V
  & 会把训练/长序列 kernel 优化和 decode 状态复用混为一谈\\
\bottomrule
\end{boundarytable}

\section{Attention 的对象：Query 提问，Key 匹配，Value 提供内容}

设 query 序列长度为 $N_q$，key/value 序列长度为 $N_k$。一个 attention head 中：

$$
\boldsymbol{Q}\in\mathbb{R}^{N_q\times d_k},\qquad
\boldsymbol{K}\in\mathbb{R}^{N_k\times d_k},\qquad
\boldsymbol{V}\in\mathbb{R}^{N_k\times d_v}.
$$

Scaled dot-product attention 定义为

$$
\boldsymbol{S}
:=\frac{\boldsymbol{Q}\boldsymbol{K}^{\mathsf T}}{\sqrt{d_k}},
$$

$$
\boldsymbol{P}
:=\operatorname{softmax}_{\text{row}}(\boldsymbol{S}+\boldsymbol{M}),
$$

$$
\boxed{
\boldsymbol{O}
=\boldsymbol{P}\boldsymbol{V}
}.
$$

这里 $\boldsymbol{M}$ 表示可选的 mask / additive bias。若某个位置不允许被关注，可以在 softmax 前把对应 score 设为 $-\infty$，使其指数权重为零。

\begin{handbookdiagram}
\centering
\begin{tikzpicture}[
  >=Stealth,
  font=\small,
  color=themefg,
  text=themefg,
  tensor/.style={draw=themecurve,fill=themecurve!11!themebg,rounded corners=4pt,minimum width=2.05cm,minimum height=0.82cm,align=center,line width=0.95pt},
  op/.style={draw=themeamber,fill=themeamber!10!themebg,rounded corners=4pt,minimum width=2.45cm,minimum height=0.88cm,align=center,line width=0.95pt},
  weight/.style={draw=themegreen,fill=themegreen!10!themebg,rounded corners=4pt,minimum width=2.25cm,minimum height=0.82cm,align=center,line width=0.95pt},
  flow/.style={->,draw=themefg!90,line width=1.05pt}
]

\node[tensor] (q) at (0,1.65) {$Q$\\\scriptsize 当前要查询什么};
\node[tensor] (k) at (0,0) {$K$\\\scriptsize 候选怎样被匹配};
\node[tensor] (v) at (7.35,-1.8) {$V$\\\scriptsize 候选提供的内容};

\node[op] (score) at (3.2,0.82) {$S=QK^{\mathsf T}/\sqrt{d_k}$\\\scriptsize pairwise score};
\node[op] (mask) at (6.15,0.82) {$S+M$\\\scriptsize 可选 mask / bias};
\node[weight] (p) at (9.0,0.82) {$P=\operatorname{softmax}_{row}(S+M)$\\\scriptsize 路由权重};
\node[op] (mix) at (11.85,0.82) {$PV$\\\scriptsize weighted aggregation};
\node[tensor] (o) at (14.35,0.82) {$O$\\\scriptsize 路由后的表示};

\draw[flow] (q.east) -- ++(0.75,0) |- (score.north west);
\draw[flow] (k.east) -- ++(0.75,0) |- (score.south west);
\draw[flow] (score) -- (mask);
\draw[flow] (mask) -- (p);
\draw[flow] (p) -- (mix);
\draw[flow] (v.east) -- ++(2.85,0) |- (mix.south);
\draw[flow] (mix) -- (o);

\node[anchor=west,font=\scriptsize,color=themegray] at (0,-2.55)
  {位置只表示数据依赖顺序；框之间的几何距离不表示计算耗时或重要性。};

\end{tikzpicture}

\diagramcaption{Scaled dot-product Attention 的信息流。箭头表示张量数据依赖；$QK^{\mathsf T}$ 产生路由分数，row-wise softmax 把每行变成归一化权重，再用权重汇总 $V$。}
\end{handbookdiagram}

$Q,K,V$ 的角色不要靠字母死记。对第 $i$ 个 query 与第 $j$ 个 candidate key：

$$
S_{ij}
=\frac{\boldsymbol{q}_i^{\mathsf T}\boldsymbol{k}_j}{\sqrt{d_k}}.
$$

$S_{ij}$ 回答“第 $i$ 个位置当前的查询与第 $j$ 个位置提供的匹配索引有多相容”。Softmax 把同一 query 对所有 key 的相对分数变成一组权重；真正被汇总的内容是 $\boldsymbol{v}_j$。

因此可以把一行输出写成

$$
\boldsymbol{o}_i
=\sum_{j=1}^{N_k}P_{ij}\boldsymbol{v}_j.
$$

这说明 Attention 是一种\textbf{内容依赖的加权路由}：权重由 $Q,K$ 决定，被路由的载荷由 $V$ 提供。

\subsection{贯穿母例：同一个 query 怎样把路由权重变成输出}

考虑一个 query 对三个 key 的缩放后 score：

$$
\boldsymbol{s}=(1,0,2),
$$

并令三个 value 为

$$
\boldsymbol{v}_1=(1,0),\qquad
\boldsymbol{v}_2=(0,1),\qquad
\boldsymbol{v}_3=(1,1).
$$

以最大 score $2$ 为指数基准，softmax 权重与

$$
\left(\mathrm{e}^{-1},\mathrm{e}^{-2},1\right)
$$

成比例，因此

$$
\boldsymbol{o}
=
\frac{
\mathrm{e}^{-1}\boldsymbol{v}_1
+\mathrm{e}^{-2}\boldsymbol{v}_2
+\boldsymbol{v}_3
}{1+\mathrm{e}^{-1}+\mathrm{e}^{-2}}
=
\frac{(1+\mathrm{e}^{-1},\ 1+\mathrm{e}^{-2})}
{1+\mathrm{e}^{-1}+\mathrm{e}^{-2}}.
$$

这个例子后面会再次出现。先把它看成普通 Attention：score 决定三条信息边的相对权重，输出只是对应 value 的加权和。进入 FlashAttention 后，数学目标不变，只改变这三个权重与加权和的计算顺序。

\subsection{为什么需要缩放因子}

若 $q_r,k_r$ 的各维分量近似独立、均值为零、方差量级相近，则点积

$$
\boldsymbol{q}^{\mathsf T}\boldsymbol{k}
=\sum_{r=1}^{d_k}q_rk_r
$$

的方差会随 $d_k$ 增长。较大的 logits 会把 softmax 推入更尖锐的区域，使梯度尺度更难控制。除以 $\sqrt{d_k}$ 的作用是让 score 的典型尺度不随 head dimension 线性扩张。

这是尺度控制，不意味着所有真实训练分布都严格满足独立同分布假设；该推导解释设计动机，而不是对任意模型激活分布的精确统计定律。

\section{Self、Cross、Causal：改变的是信息可达关系}

\subsection{Self-Attention}

Self-attention 中 $Q,K,V$ 都来自同一组输入表示的不同投影。于是每个位置可以根据当前内容重新读取同一序列中的其他位置。

\subsection{Cross-Attention}

Cross-attention 中 query 与 key/value 来自不同表示集合，例如 decoder state 查询 encoder output。此时一般有

$$
N_q\neq N_k,
$$

所以输出 shape 跟 query 侧一致：

$$
\boldsymbol{O}\in\mathbb{R}^{N_q\times d_v}.
$$

\subsection{Causal Attention}

自回归场景要求位置 $i$ 不能读取未来位置 $j>i$。定义

$$
M_{ij}
=
\begin{cases}
0,& j\le i,\\
-\infty,& j>i.
\end{cases}
$$

于是

$$
P_{ij}=0\qquad(j>i).
$$

Causal mask 改变的是\textbf{允许的信息边}，不是 softmax 本身的数学定义。

\section{Multi-Head Attention：并行建立多个路由子空间}

若模型维度为 $d_{\text{model}}$，多头注意力通过不同参数投影得到多组

$$
\boldsymbol{Q}^{(h)},\boldsymbol{K}^{(h)},\boldsymbol{V}^{(h)},
$$

分别执行 attention：

$$
\operatorname{head}_h
=\operatorname{Attention}
\left(
\boldsymbol{Q}^{(h)},
\boldsymbol{K}^{(h)},
\boldsymbol{V}^{(h)}
\right).
$$

然后拼接并线性投影：

$$
\operatorname{MHA}(X)
=\operatorname{Concat}(\operatorname{head}_1,\ldots,\operatorname{head}_H)\boldsymbol{W}_O.
$$

多头并不保证“每个 head 必然学一种人类可解释关系”。更准确地说，模型获得多组独立参数化的相似度与 value 投影通道，可以并行形成不同的信息路由模式。

\section{位置信息：纯内容路由本身不知道顺序从哪里来}

如果同时对输入位置做同一个置换，而没有额外 positional signal，self-attention 会相应置换输出；它本身不会凭空知道“第 3 个 token”和“第 30 个 token”的绝对或相对顺序。

因此，如果任务不仅需要 causal mask 这类“谁能看谁”的可达性约束，还需要表达 token 的绝对位置、相对距离或顺序关系，就要通过某种 positional mechanism 把这些信息注入表示或 attention score。具体实现可以是绝对位置编码、相对位置 bias、旋转位置表示等；这些实现会演化，但它们都在补充纯内容相似度本身没有提供的位置关系：

$$
\boxed{
\text{content identity}
+\text{position relation}
\longrightarrow
\text{order-aware routing}
}.
$$

这里的箭头表示信息组合，不表示数学蕴含。

\section{Transformer block：把路由与逐位置变换交替堆叠}

Attention 负责跨位置交换信息，feed-forward network（FFN）负责对每个位置的 channel representation 做非线性变换。Residual connection 为原表示提供直接通路；normalization 控制子层输入或输出的尺度，并改变深层网络的优化动态。具体效果取决于 norm 类型与放置位置，因此这里只保留架构职责，不把某一种 placement 当成唯一标准。

现代 Transformer 存在 Pre-Norm、Post-Norm、RMSNorm 等多种变体，因此不能把某一种 norm placement 写成“Transformer 的定义”。下面只用常见 Pre-Norm 形式展示信息流：

\begin{handbookdiagram}
\centering
\begin{tikzpicture}[
  >=Stealth,
  font=\small,
  color=themefg,
  text=themefg,
  state/.style={draw=themecurve,fill=themecurve!11!themebg,rounded corners=4pt,minimum width=1.9cm,minimum height=0.82cm,align=center,line width=0.95pt},
  op/.style={draw=themeamber,fill=themeamber!10!themebg,rounded corners=4pt,minimum width=2.1cm,minimum height=0.86cm,align=center,line width=0.95pt},
  add/.style={circle,draw=themegreen,fill=themegreen!10!themebg,minimum size=0.68cm,line width=1pt,font=\bfseries},
  flow/.style={->,draw=themefg!90,line width=1.05pt},
  residual/.style={->,draw=themecurve,line width=1pt,dashed}
]

\node[state] (x) at (0,0) {$X$};
\node[op] (n1) at (2.35,0) {Norm};
\node[op] (attn) at (5.0,0) {Multi-Head\\Attention};
\node[add] (a1) at (7.25,0) {$+$};
\node[op] (n2) at (9.45,0) {Norm};
\node[op] (ffn) at (12.0,0) {FFN};
\node[add] (a2) at (14.2,0) {$+$};
\node[state] (y) at (16.25,0) {$Y$};

\draw[flow] (x) -- (n1);
\draw[flow] (n1) -- (attn);
\draw[flow] (attn) -- (a1);
\draw[flow] (a1) -- (n2);
\draw[flow] (n2) -- (ffn);
\draw[flow] (ffn) -- (a2);
\draw[flow] (a2) -- (y);

\draw[residual] (x.north) -- ++(0,1.35) -| (a1.north);
\draw[residual] (a1.south) -- ++(0,-1.35) -| (a2.south);

\node[anchor=west,font=\scriptsize,color=themecurve] at (2.2,1.55) {residual：保留输入并与子层输出相加};
\node[anchor=west,font=\scriptsize,color=themegray] at (0,-1.9)
  {此图画的是常见 Pre-Norm 结构；Post-Norm 等变体改变 Norm 位置，但不改变 Attention 与 FFN 的职责边界。};

\end{tikzpicture}

\diagramcaption{一个常见 Pre-Norm Transformer block 的信息流。实线箭头表示前向数据流，绕行边表示 residual addition；其他 norm placement 仍属于同一 architecture family。}
\end{handbookdiagram}

抽象地看，一层不断交替执行两件事：

\begin{enumerate}
  \item \textbf{跨位置路由}：当前位置从其他位置读取什么；
  \item \textbf{逐位置计算}：读取后的表示在 feature/channel 维度上怎样变换。
\end{enumerate}

所以 Transformer 不是“Attention 连续堆叠”。没有 FFN、residual、normalization 与 position，不能把单独 attention 算子直接等同完整 Transformer block。

\section{Attention 扩展的两个轴：序列长度与模型深度不能混为一谈}

到这里需要先把一个容易被新模型术语混淆的边界钉死：\textbf{“让更长序列算得更便宜”和“让更深网络选择更早层表示”是两个不同问题。} 它们都可能使用 attention / routing 语言，但路由发生的轴不同。

\begin{handbookdiagram}
\centering
\begin{tikzpicture}[
  >=Stealth,
  font=\small,
  color=themefg,
  text=themefg,
  base/.style={draw=themecurve,fill=themecurve!11!themebg,rounded corners=5pt,minimum width=3.0cm,minimum height=0.92cm,align=center,line width=1pt},
  seq/.style={draw=themeamber,fill=themeamber!10!themebg,rounded corners=5pt,minimum width=3.15cm,minimum height=0.92cm,align=center,line width=1pt},
  depth/.style={draw=themegreen,fill=themegreen!10!themebg,rounded corners=5pt,minimum width=3.15cm,minimum height=0.92cm,align=center,line width=1pt},
  impl/.style={draw=themealert,fill=themealert!8!themebg,rounded corners=5pt,minimum width=3.15cm,minimum height=0.92cm,align=center,line width=1pt},
  flow/.style={->,draw=themefg!88,line width=1.0pt},
  combine/.style={->,draw=themeamber,line width=0.98pt},
  relation/.style={->,draw=themegray,dashed,line width=0.95pt},
  implrel/.style={<->,draw=themealert,dashed,line width=0.95pt}
]

% Sequence-axis architecture choices.
\begin{pgfonlayer}{background}
  \fill[fill=themeamber!4!themebg,draw=themeamber!40,dashed,rounded corners=7pt]
    (-0.1,2.0) rectangle (15.2,6.2);
  \fill[fill=themegreen!4!themebg,draw=themegreen!40,dashed,rounded corners=7pt]
    (-0.1,-2.2) rectangle (15.2,1.45);
\end{pgfonlayer}

\node[anchor=west,font=\bfseries,color=themeamber] at (0.15,5.8) {Sequence axis：token 之间怎样交换信息};
\node[base] (full) at (3.0,4.35) {Full Softmax Attention\\\scriptsize explicit token-to-token interaction};
\node[seq] (linear) at (8.0,4.35) {Recurrent / Linear Attention\\\scriptsize finite-state sequence mixer};
\node[seq] (hybrid) at (13.0,4.35) {Hybrid Attention\\\scriptsize recurrent + periodic global mixing};

\draw[combine] (full.north) -- ++(0,0.58) -| node[pos=0.62,above=1pt,font=\scriptsize,color=themeamber]{保留 global path} (hybrid.north);
\draw[combine] (linear.east) -- node[above=1pt,font=\scriptsize,color=themeamber]{组合} (hybrid.west);

% FlashAttention is an implementation relation to full Softmax, not a sibling architecture choice.
\node[impl] (flash) at (3.0,2.65) {FlashAttention\\\scriptsize same mathematical mapping · tiled I/O};
\draw[implrel] (full.south) -- node[right=2pt,font=\scriptsize,color=themealert]{数学映射相同} (flash.north);

\node[seq,minimum width=3.75cm] (state) at (8.0,2.65) {$S_t\leftarrow F(S_{t-1},k_t,v_t)$\\$o_t=G(S_t,q_t)$};
\draw[relation] (linear.south) -- node[right=2pt,font=\scriptsize,color=themegray]{典型状态化形式} (state.north);

% Depth-axis routing is separate from sequence mixing.
\node[anchor=west,font=\bfseries,color=themegreen] at (0.15,1.05) {Depth axis：当前层怎样读取更早层表示};
\node[base] (resid) at (3.0,-0.35) {Standard Residual\\$h_{\ell+1}=h_\ell+f_\ell(h_\ell)$};
\node[depth] (attnres) at (8.0,-0.35) {Attention Residuals\\\scriptsize learned routing over earlier layers};
\node[depth] (variants) at (13.0,-0.35) {Block / Delta AttnRes\\\scriptsize compress sources / change routed values};
\draw[flow,draw=themegreen] (resid) -- node[above=1pt,font=\scriptsize,color=themegreen]{显式 depth routing} (attnres);
\draw[flow,draw=themegreen] (attnres) -- node[above=1pt,font=\scriptsize,color=themegreen]{继续改写 source} (variants);

\node[anchor=west,font=\scriptsize,color=themegray] at (0.15,-1.55)
  {关键分层：FlashAttention 只改变 Full Softmax 的执行方式；Linear/Hybrid 改写 sequence mixer；AttnRes 则在另一条 depth 轴上路由更早层表示。};

\end{tikzpicture}

\diagramcaption{Attention 相关扩展的两个架构轴与一个实现层。上半部分沿 sequence 轴区分 Full Softmax、recurrent/linear 与 hybrid mixer；FlashAttention 单独挂在 Full Softmax 下方，表示“数学映射相同、执行方式不同”。下半部分沿 depth 轴讨论 residual / AttnRes 对更早层表示的聚合。}
\end{handbookdiagram}

\subsection{FlashAttention 与 Linear Attention：一个优化实现，一个改变映射}

标准 causal Softmax Attention 的第 $t$ 个输出会显式依赖此前所有 key/value：

$$
\boldsymbol{o}_t
=\sum_{j\le t}
\operatorname{softmax}(\boldsymbol{s}_t)_j\boldsymbol{v}_j.
$$

FlashAttention 仍然计算这一个映射，只是通过 tiling、online softmax 和片上复用避免物化完整 $S/P$。因此它仍然属于\textbf{精确 Softmax Attention 的 I/O-aware 实现}，核心序列计算没有从二次 all-pairs interaction 变成线性 recurrence。

Recurrent / linear attention 走的是另一条路：不再让当前 query 与全部历史 key 显式形成完整 score row，而是把过去压缩进固定大小的 recurrent state。抽象形式可以写成

$$
\boxed{
S_t=F(S_{t-1},k_t,v_t),
\qquad
 o_t=G(S_t,q_t)
}.
$$

这样每一步读取的是状态 $S_t$，而不是长度随历史增长的完整 KV 集合。代价是：\textbf{它已经改变了 token mixing 的数学结构}。因此它不能被描述为“更快的 FlashAttention”，也不能默认与标准 Softmax Attention 数值等价。

\subsection{Kimi Delta Attention：有限状态记忆中的衰减、纠正与写入}

Kimi Delta Attention（KDA）是当前 recurrent linear attention 的一个具体例子。它从 Gated DeltaNet 发展而来，用 fine-grained decay gate 控制固定大小的矩阵状态。Moonshot 当前公开的 KDA forward 将每个 head 的状态写成

$$
\boxed{
S_t
=
\left(I-\beta_t k_tk_t^{\mathsf T}\right)
\operatorname{Diag}(\alpha_t)S_{t-1}
+\beta_t k_tv_t^{\mathsf T}
}
$$

并读取

$$
\boxed{o_t=S_t^{\mathsf T}q_t}.
$$

这里 $S_t\in\mathbb{R}^{d_k\times d_v}$ 的 shape 不随序列长度增长；$\alpha_t\in(0,1)^{d_k}$ 是逐 channel 的 retention / decay gate，$\beta_t$ 控制当前 delta-rule correction / write 的强度。把更新拆开看更容易理解：

\begin{enumerate}
  \item 先用 $\alpha_t$ 对旧状态做选择性衰减；
  \item 再沿当前 key 方向修正旧关联，避免简单叠加不断污染状态；
  \item 最后写入当前 $k_t v_t^{\mathsf T}$ 关联；
  \item query $q_t$ 从更新后的有限状态中读取输出。
\end{enumerate}

因此把 KDA 说成“滚动平均”会丢掉最关键的结构：它维护的是 key--value associative matrix state，并带有衰减、delta correction 与新关联写入，而不是维护一个标量均值。

有限状态也带来新的 trade-off。历史被压缩后，系统不再拥有标准 Softmax Attention 那种对任意过去 token 的显式独立访问；需要通过 state update 学会保留、覆盖和读取有用关联。换来的则是更稳定的长序列状态规模与更低的 decode cache 压力。

\subsection{从 Gated DeltaNet 到 KDA 再到 Gated DeltaNet-2：有限状态记忆到底要控制什么}

把近期 linear-attention 进展放在一起看，真正稳定的研究问题不是“谁把 $O(N^2)$ 变成 $O(N)$”，而是：\textbf{当全部历史被压缩进有限状态以后，模型怎样精确决定保留、擦除和写入什么？}

可以把一个 recurrent memory update 抽象成三个相互关联但不完全相同的控制量：

\begin{enumerate}
  \item \textbf{Retention / Decay}：旧状态哪些 channel 继续保留，哪些应该快速衰减；
  \item \textbf{Erase / Correction}：当前 key 指向的旧关联是否需要先被纠正或擦除；
  \item \textbf{Write}：新的 key--value 关联以多大强度写入状态。
\end{enumerate}

Gated DeltaNet 把 adaptive forgetting 与 delta-rule correction 结合起来；KDA 的关键推进是把 decay 从更粗粒度控制扩展为\textbf{逐 channel 的 fine-grained decay}，从而更细致地管理有限状态容量。2026 年的 Gated DeltaNet-2 又进一步指出：erase 与 write 本身仍然不必绑定在同一个门控强度上，于是把两者解耦为独立控制。

因此这条演进最值得保留的不是三个模型名，而是一个更一般的有限记忆设计原则：

$$
\boxed{
\text{compressed recurrent memory}
\Rightarrow
\text{forget / erase / write 是需要分别建模的控制问题}
}
$$

这不是“所有 linear attention 都必须采用三个独立 gate”的定理；它表达的是设计诊断：若这些作用长期被同一个粗粒度控制量绑定，模型可能难以同时处理过度遗忘、旧关联污染和新信息写入不足。

Gated DeltaNet-2 的当前结果只来自作者报告的有限模型规模与训练设置；它说明“解耦 erase/write”是有实验证据的研究方向，但还不能据此宣布它已经取代 KDA 或 full attention。

\subsection{Hybrid Attention：Linear 与 Global Attention 可以共存}

Linear attention 与 full/global attention 也不是必须二选一。Kimi Linear 的公开架构采用 3:1 的 KDA-to-global MLA 层比例：多数层利用 recurrent state 提高长上下文效率，周期性 global attention 层保留更直接的全局 token interaction。

这给出一个可复用的设计视角：

$$
\boxed{
\text{cheap recurrent mixing}
+\text{occasional global mixing}
}
$$

是否值得这样混合，取决于质量、长上下文需求、KV-cache 成本、训练并行化和目标硬件；“linear”并不意味着所有 full attention 都应被删除。

\subsection{Attention Residuals：路由发生在 depth，而不是 token sequence}

标准 Pre-Norm residual stack 会把各层产生的更新沿深度方向持续累加。展开多层以后，可以把它理解为更早层贡献以固定加法权重进入后续 hidden state。Attention Residuals（AttnRes）改变的是这一\textbf{跨层聚合规则}：当前层不再只接受固定单位权重的累计结果，而是通过 learned softmax attention 在多个较早 layer outputs 之间选择性聚合。

因此这里的“attention”回答的是：

> 当前第 $\ell$ 层应该从哪些较早深度的表示读取多少信息？

而普通 self-attention 回答的是：

> 当前 token 应该从哪些 token positions 读取多少信息？

把 Full AttnRes 抽象化，可以写成：

$$
\pi_{\ell,j}
=\operatorname{softmax}_j(s_{\ell,j}),
\qquad
r_\ell
=\sum_{j<\ell}\pi_{\ell,j}h_j,
$$

其中 $h_j$ 是更早深度的候选表示，$s_{\ell,j}$ 是当前层对这些 depth sources 的 learned score。标准 residual 相当于把历史贡献通过固定加法不断累积；AttnRes 则把“哪些早期表示值得继续进入当前层”变成 input-dependent routing。

两者一个沿 \textbf{depth axis} 路由，一个沿 \textbf{sequence axis} 路由。Full AttnRes 若直接保存并访问所有早期 layer outputs，会带来随深度增长的 activation memory 与 pipeline communication。Block AttnRes 因此把若干相邻层压成 block-level sources：当前层只在少量 block representations 之间做 depth-wise routing，换取更低的状态与通信开销。

\subsection{从 Block AttnRes 到 Delta AttnRes：下一步研究的是“应该路由哪种深度表示”}

AttnRes 把“跨层路由”显式化以后，又暴露出新的设计自由度：被路由的 value 到底应该是累计 hidden state，还是每一层真正新增的信息？

2026 年后续工作 Delta Attention Residuals 提出使用

$$
\boxed{
\Delta h_j=h_{j+1}-h_j
}
$$

作为 depth-routing 的候选表示，而不是直接在高度相关的累计 $h_j$ 之间选择。作者的解释是：深层 cumulative states 往往彼此高度冗余，使 softmax routing 越来越接近低对比度；layer delta 更直接表示“这一层新写入了什么”。其在作者报告的 $220$M--$7.6$B 模型上优于标准 residual 与原始 AttnRes。

这里必须保持证据边界：Delta AttnRes 是独立后续研究，不是 Kimi K3 的已发布结构；当前结果说明“\textbf{depth routing 的 source representation 选择本身是一阶设计问题}”，而不是证明所有大模型都应把 AttnRes 改成 delta routing。

\begin{boundarytable}
\boundarytablehead
FlashAttention & $\neq$ & KDA / Linear Attention
  & 前者保持标准 Softmax Attention 映射并优化 I/O；后者用 recurrent finite state 改写 sequence mixer
  & 会把 exact kernel optimization 与 sequence-mixer replacement 混成一件事\\
Token Attention & $\neq$ & Attention Residuals
  & 前者沿 sequence positions 路由；后者沿 model depth 在较早 layer representations 之间路由
  & 会误以为 AttnRes 是另一种长上下文 token attention\\
MoE Routing & $\neq$ & Attention Routing
  & MoE 决定一个 token 进入哪些 expert/FFN 参数子网络；Attention 决定表示从哪些 token 或 layer source 聚合信息
  & 会把“只激活少量 experts”误解成 attention 变稀疏\\
\bottomrule
\end{boundarytable}

\subsection{Kimi K3：同一个大模型同时优化 sequence、depth 与 channel 三条信息流}

Kimi K3 值得作为架构阅读案例，不是因为“2.8T 参数很大”，而是因为它把三类本来容易混在一起的路由问题同时放进一个模型：

\begin{handbooktable}{L{0.22\linewidth}L{0.28\linewidth}Y}
\toprule
\textbf{轴} & \textbf{K3 中的机制} & \textbf{真正解决的问题}\\
\midrule
Sequence mixing & 69 KDA + 24 Gated MLA layers & 历史 token 怎样以有限状态高效压缩，同时周期性保留 global content interaction\\
Depth mixing & Block Attention Residuals & 当前层怎样选择性读取较早 depth representations，缓解固定 residual accumulation 的 dilution\\
Channel / parameter mixing & Stable LatentMoE，896 routed experts 中每 token 选择 16 个，另有 shared experts & 怎样把参数容量扩得很大，但只让少量 expert 子网络参与当前 token 计算\\
\bottomrule
\end{handbooktable}

\begin{handbookdiagram}
\centering
\begin{tikzpicture}[
  >=Stealth,
  font=\small,
  color=themefg,
  text=themefg,
  seq/.style={draw=themeamber,fill=themeamber!10!themebg,rounded corners=5pt,minimum width=3.55cm,minimum height=0.94cm,align=center,line width=1pt},
  depth/.style={draw=themegreen,fill=themegreen!10!themebg,rounded corners=5pt,minimum width=3.55cm,minimum height=0.94cm,align=center,line width=1pt},
  channel/.style={draw=themecurve,fill=themecurve!11!themebg,rounded corners=5pt,minimum width=3.55cm,minimum height=0.94cm,align=center,line width=1pt},
  snapshot/.style={draw=themegray,fill=themegray!8!themebg,rounded corners=5pt,minimum width=6.1cm,minimum height=0.82cm,align=center,line width=0.9pt},
  flow/.style={->,draw=themefg!88,line width=1pt}
]

% Column semantics are shared by all three rows.
\node[anchor=west,font=\scriptsize\bfseries,color=themegray] at (0.25,5.65) {信息来源};
\node[anchor=west,font=\scriptsize\bfseries,color=themegray] at (6.55,5.65) {路由 / 混合机制};
\node[anchor=west,font=\scriptsize\bfseries,color=themegray] at (12.55,5.65) {送往哪里};

% Sequence axis.
\node[anchor=west,font=\bfseries,color=themeamber] at (0.25,4.85) {Sequence axis};
\node[seq] (seqsrc) at (2.65,3.75) {History tokens / state\\\scriptsize full KV interaction or recurrent memory};
\node[seq] (seqmix) at (8.05,3.75) {69 KDA + 24 Gated MLA\\\scriptsize recurrent mixing + periodic global mixing};
\node[seq] (seqdst) at (13.45,3.75) {Current token context\\\scriptsize sequence-wise information aggregation};
\draw[flow,draw=themeamber] (seqsrc) -- (seqmix);
\draw[flow,draw=themeamber] (seqmix) -- (seqdst);

% Depth axis.
\node[anchor=west,font=\bfseries,color=themegreen] at (0.25,2.65) {Depth axis};
\node[depth] (depsrc) at (2.65,1.55) {Earlier block representations\\\scriptsize candidate depth sources};
\node[depth] (depmix) at (8.05,1.55) {Block AttnRes\\\scriptsize learned softmax across depth sources};
\node[depth] (depdst) at (13.45,1.55) {Current layer input\\\scriptsize selective cross-layer aggregation};
\draw[flow,draw=themegreen] (depsrc) -- (depmix);
\draw[flow,draw=themegreen] (depmix) -- (depdst);

% Channel / parameter axis.
\node[anchor=west,font=\bfseries,color=themecurve] at (0.25,0.45) {Channel / parameter axis};
\node[channel] (chsrc) at (2.65,-0.65) {Hidden state $x$\\\scriptsize residual width $7168$};
\node[channel] (chmix) at (8.05,-0.65) {Stable LatentMoE\\\scriptsize down-project $3584$ · Top-16 / 896 experts};
\node[channel] (chdst) at (13.45,-0.65) {Expert contribution\\\scriptsize up-project to residual width};
\draw[flow,draw=themecurve] (chsrc) -- (chmix);
\draw[flow,draw=themecurve] (chmix) -- (chdst);

\node[snapshot] at (8.05,-2.5)
  {Kimi K3 · 2026 architecture snapshot：2.8T total · 104B activated · 93 layers · $1{,}048{,}576$ context};

\node[anchor=west,font=\scriptsize,color=themegray] at (0.25,-3.45)
  {三行共享同一读法：先找“从哪里读”，再看“怎样路由/混合”，最后看“写到哪里”。三个轴可以同时存在，但路由对象完全不同。};

\end{tikzpicture}

\diagramcaption{Kimi K3 的三轴分解。Sequence mixing、depth mixing 与 expert/channel mixing 是三个正交程度较高的架构问题；同一 block 可以同时包含这些机制，但一个轴的“稀疏/路由”不能拿来解释另一个轴。}
\end{handbookdiagram}

当前官方配置是 93 layers、2.8T total parameters、约 104B activated parameters、$1{,}048{,}576$ context length。69 个 KDA 与 24 个 Gated MLA 形成约 3:1 的 sequence-mixing hybrid；这比“全部改成 linear attention”更能说明当前工程取向：\textbf{用多数便宜 recurrent layers 承担长序列状态更新，再周期性插入 global layers 补充显式全局交互。}

Stable LatentMoE 又解决完全不同的 channel/parameter 问题。按 K3 当前公开结构，可以把 routed branch 压成：

$$
z=W_{\mathrm{down}}x\in\mathbb{R}^{3584},
$$

$$
u
=\sum_{i\in\mathcal{T}(x)}p_iE_i(z),
\qquad |\mathcal{T}(x)|=16,
$$

$$
y_{\mathrm{routed}}
=W_{\mathrm{up}}\operatorname{RMSNorm}(u)
\in\mathbb{R}^{7168}.
$$

Router 仍根据 full-width representation 选择 896 routed experts 中的 Top-16，但真正 routed-expert computation 在较窄的 latent width $3584$ 中进行，再投影回 residual width $7168$。这样做的系统意义是：\textbf{参数容量可以通过大量 experts 扩展，而每 token 的 expert weight traffic、expert compute 与 distributed dispatch payload 不必按 full hidden width 同比例增长。}

在高稀疏 MoE 中，另一个一阶问题是负载均衡：若少数 experts 长期过热，理论稀疏性并不会自动转化为训练吞吐。K3 的 Stable LatentMoE 使用 Quantile Balancing，根据 router-score quantiles 调整 expert allocation bias；其目的不是改变最终 semantic routing 定义，而是让极高 expert sparsity 下的训练负载更稳定。K3 技术报告所称相对 K2 约 $2.5\times$ 的 overall scaling-efficiency improvement 是\textbf{整套架构、训练和数据 recipe 的作者报告结果}，不能单独归因于 KDA、AttnRes 或 LatentMoE 中某一个模块。

因此，读类似新模型时最稳的拆法不是背产品参数，而是依次问：

\begin{enumerate}
  \item sequence 方向的信息如何混合；
  \item depth 方向的信息如何保留与选择；
  \item channel / parameter 容量如何激活；
  \item 这些数学结构怎样映射到 kernel、通信、cache 与并行系统。
\end{enumerate}

MoE 的一般路由、负载均衡与分布式 expert parallel 不是本 Attention Topic 的 Canonical Owner，这里只保留理解 K3 组合架构所需的接口。

\section{长序列的第二个问题：数学正确不等于硬件高效}

对 self-attention，若

$$
\boldsymbol{Q},\boldsymbol{K},\boldsymbol{V}\in\mathbb{R}^{N\times d},
$$

则

$$
\boldsymbol{S}=\boldsymbol{Q}\boldsymbol{K}^{\mathsf T}
\in\mathbb{R}^{N\times N},
$$

以及

$$
\boldsymbol{P}=\operatorname{softmax}(\boldsymbol{S})
\in\mathbb{R}^{N\times N}.
$$

标准 attention 的核心计算仍具有二次序列规模：

$$
\Theta(N^2d).
$$

朴素 kernel chain 还可能把 $S$、$P$ 分别写入 device global memory；在 HBM-backed GPU 上，这对应把大型中间矩阵写入 HBM，再由后续 kernel 重新读回：

\begin{processchain}
  \processstage{$Q,K$}
  \processstage{GEMM 得 $S$}
  \processstage{$S$ 写 HBM}
  \processstage{Softmax 得 $P$}
  \processstage{$P$ 写 HBM}
  \processstage{$PV$ 得 $O$}
\end{processchain}

因此运行时间不仅由 FLOPs 决定，还受 data movement 限制。可以用 arithmetic intensity 建立最小性能语言：

$$
\operatorname{AI}
:=\frac{\text{FLOPs}}{\text{Bytes moved}}.
$$

当 memory traffic 相对于计算量过大时，增加更多算术单元并不会自动线性加速 kernel。

\begin{mechanism}{FlashAttention 的根本转换}
FlashAttention 不改变标准 Attention 的数学映射，也没有把 $\Theta(N^2)$ attention 工作变成线性计算。它改变的是执行顺序：把 $Q,K,V$ 分块，在片上尽可能完成 score、softmax 状态更新和 value accumulation，避免把完整 $N\times N$ 的 $S$、$P$ 反复物化到 device global memory。浮点运算次序可能与朴素实现不同，因此这里的“相同”指数学语义相同，而不是逐 bit 相同。
\end{mechanism}

\begin{handbookdiagram}
\centering
\begin{tikzpicture}[
  >=Stealth,
  font=\small,
  color=themefg,
  text=themefg,
  data/.style={draw=themecurve,fill=themecurve!11!themebg,rounded corners=4pt,minimum width=1.9cm,minimum height=0.72cm,align=center,line width=0.9pt},
  op/.style={draw=themeamber,fill=themeamber!10!themebg,rounded corners=4pt,minimum width=2.35cm,minimum height=0.78cm,align=center,line width=0.9pt},
  mem/.style={draw=themealert,fill=themealert!9!themebg,rounded corners=4pt,minimum width=2.25cm,minimum height=0.76cm,align=center,line width=0.9pt},
  state/.style={draw=themegreen,fill=themegreen!10!themebg,rounded corners=4pt,minimum width=2.7cm,minimum height=0.86cm,align=center,line width=0.95pt},
  flow/.style={->,draw=themefg!90,line width=1pt},
  stream/.style={->,draw=themegreen,line width=0.98pt},
  loop/.style={->,draw=themeamber,dashed,line width=0.98pt}
]

\begin{pgfonlayer}{background}
  \fill[fill=themealert!4!themebg,draw=themealert!40,dashed,rounded corners=7pt]
    (-0.4,-4.65) rectangle (7.0,4.6);
  \fill[fill=themegreen!4!themebg,draw=themegreen!40,dashed,rounded corners=7pt]
    (7.5,-4.65) rectangle (16.3,4.6);
  \fill[fill=themealert!3!themebg,draw=themealert!35,densely dashed,rounded corners=5pt]
    (7.85,-3.95) rectangle (10.25,3.7);
\end{pgfonlayer}

\node[anchor=west,font=\bfseries,color=themealert] at (-0.05,4.18) {朴素多 Kernel：完整中间矩阵落回 HBM};
\node[anchor=west,font=\bfseries,color=themegreen] at (7.85,4.18) {Tiled exact Attention：HBM 流式供给，片上维护状态};

% Naive path: each large intermediate is materialized between kernels.
\node[data] (qk) at (1.15,2.85) {$Q,K$};
\node[op] (gemm) at (4.2,2.85) {$S=QK^{\mathsf T}$};
\node[mem] (hbmS) at (4.2,1.35) {HBM\\\scriptsize materialize $S$};
\node[op] (soft) at (4.2,-0.15) {row Softmax};
\node[mem] (hbmP) at (4.2,-1.65) {HBM\\\scriptsize materialize $P$};
\node[data] (v) at (1.15,-3.05) {$V$};
\node[op] (pv) at (4.2,-3.05) {$O=PV$};
\node[data] (o1) at (6.05,-3.05) {$O$};
\draw[flow] (qk) -- (gemm);
\draw[flow] (gemm) -- (hbmS);
\draw[flow] (hbmS) -- (soft);
\draw[flow] (soft) -- (hbmP);
\draw[flow] (hbmP) -- (pv);
\draw[flow] (v) -- (pv);
\draw[flow] (pv) -- (o1);

% Right panel: HBM column on the left, on-chip working set on the right.
\node[anchor=west,font=\scriptsize\bfseries,color=themealert] at (8.15,3.35) {HBM};
\node[mem] (hq) at (9.05,2.45) {$Q_i$ source};
\node[mem] (hkv) at (9.05,0.55) {$K_j,V_j$ source\\\scriptsize streamed tile by tile};
\node[mem] (ho) at (9.05,-3.0) {$O_i$ write-back};

\node[anchor=west,font=\scriptsize\bfseries,color=themegreen] at (10.75,3.35) {on-chip working set};
\node[data] (qtile) at (11.7,2.45) {$Q_i$ resident};
\node[data] (kvtile) at (11.7,0.55) {current $(K_j,V_j)$};
\node[op] (score) at (14.55,2.05) {$S_{ij}=Q_iK_j^{\mathsf T}$\\\scriptsize ephemeral score tile};
\node[state] (online) at (14.55,0.15) {online state\\$m,l,\widetilde O$};
\node[op] (norm) at (14.55,-2.05) {$O_i=\widetilde O/l$};

\draw[stream] (hq.east) -- (qtile.west);
\draw[stream] (hkv.east) -- (kvtile.west);
\draw[flow] (qtile.east) -- ++(0.55,0) |- (score.north west);
\draw[flow] (kvtile.east) -- ++(0.55,0) |- (score.south west);
\draw[flow] (score) -- (online);
\draw[flow] (kvtile.east) -- ++(0.85,0) |- node[pos=0.72,below=1pt,font=\scriptsize]{提供 $V_j$} (online.west);
\draw[loop] (online.west) -- ++(-3.45,0) |- node[pos=0.58,left=2pt,font=\scriptsize,color=themeamber]{还有 KV tile：$j\!\leftarrow\!j+1$} (hkv.south);
\draw[flow] (online) -- (norm);
\draw[stream] (norm.west) -- ++(-2.0,0) |- (ho.east);

\node[anchor=west,font=\scriptsize,color=themegray] at (7.85,-4.28)
  {右侧循环的持久对象只有 $Q_i$、当前 KV tile 与行级 online state；$S_{ij}$ / 局部概率只在片上短暂存在，不组成完整 $N\times N$ HBM 张量。};

\end{tikzpicture}

\diagramcaption{朴素 Attention 与 tiled exact Attention 的数据流对比。左侧多次把完整中间矩阵写入 HBM；右侧固定一个 query tile，流式读取 KV tiles，并只维护行级统计量与输出累加器。图只表达数据流与驻留层次，不表示真实延迟比例。}
\end{handbookdiagram}

\section{Stable Softmax：分块之前先解决数值稳定性}

对一行 logits $\boldsymbol{x}$，softmax 为

$$
p_i
=\frac{\mathrm{e}^{x_i}}
{\sum_j\mathrm{e}^{x_j}}.
$$

直接指数化可能 overflow。利用 softmax 的平移不变性，令

$$
m:=\max_j x_j,
$$

则

$$
\boxed{
 p_i
 =\frac{\mathrm{e}^{x_i-m}}
 {\sum_j\mathrm{e}^{x_j-m}}
}.
$$

此时

$$
x_i-m\le0,
$$

所以指数项不超过 $1$。减最大值改变中间表示，却不改变最终归一化概率。

\section{Online Softmax：把“必须看完整行”改写成可合并状态}

普通 stable softmax 似乎要求先看到整行，才能知道全局最大值。若 attention score 被拆成多个 tile，这会成为关键障碍。

先看一维流式版本。处理到前 $t$ 个元素时，定义

$$
m_t
:=\max_{1\le j\le t}x_j,
$$

$$
l_t
:=\sum_{j=1}^{t}\mathrm{e}^{x_j-m_t}.
$$

加入新元素 $x_{t+1}$ 后：

$$
m_{t+1}=\max(m_t,x_{t+1}).
$$

旧的指数和原本以 $m_t$ 为基准，必须换到新的基准 $m_{t+1}$。因为

$$
\mathrm{e}^{x_j-m_{t+1}}
=\mathrm{e}^{x_j-m_t}
 \mathrm{e}^{m_t-m_{t+1}},
$$

所以

$$
\boxed{
 l_{t+1}
 =l_t\mathrm{e}^{m_t-m_{t+1}}
 +\mathrm{e}^{x_{t+1}-m_{t+1}}
}.
$$

这条递推的关键不只是“少一次遍历”，而是 $m,l$ 已经包含继续合并后续数据所需的全部行级信息。新数据到来时，不需要保存前面所有 logits。

\section{Blockwise Online Softmax：FlashAttention 的状态不变量}

固定一个 query tile $\boldsymbol{Q}_i$，依次读入 key/value tile $(\boldsymbol{K}_j,\boldsymbol{V}_j)$。当前 score tile：

$$
\boldsymbol{S}_j
=\frac{\boldsymbol{Q}_i\boldsymbol{K}_j^{\mathsf T}}{\sqrt d}.
$$

对 query tile 中每一行维护三个状态：

$$
\boldsymbol{m},\qquad
\boldsymbol{l},\qquad
\widetilde{\boldsymbol{O}}.
$$

它们的精确定义是：对目前已经处理过的 key 集合 $\mathcal{J}$，第 $r$ 行满足

$$
m_r
=\max_{j\in\mathcal{J}}S_{rj},
$$

$$
l_r
=\sum_{j\in\mathcal{J}}\mathrm{e}^{S_{rj}-m_r},
$$

$$
\widetilde{\boldsymbol{o}}_r
=\sum_{j\in\mathcal{J}}
\mathrm{e}^{S_{rj}-m_r}\boldsymbol{v}_j.
$$

因此对已经处理的部分，当前精确归一化输出就是

$$
\boldsymbol{o}_r
=\frac{\widetilde{\boldsymbol{o}}_r}{l_r}.
$$

这个定义就是 FlashAttention forward 的核心不变量。

读到新 tile 后，先计算

$$
\boldsymbol{m}_{\text{new}}
=\max\left(
\boldsymbol{m}_{\text{old}},
\operatorname{rowmax}(\boldsymbol{S}_j)
\right),
$$

以及旧状态的缩放系数

$$
\boldsymbol{\alpha}
=\mathrm{e}^{\boldsymbol{m}_{\text{old}}-\boldsymbol{m}_{\text{new}}}.
$$

当前未归一化权重为

$$
\boldsymbol{P}_j^{*}
=\mathrm{e}^{
\boldsymbol{S}_j-
\boldsymbol{m}_{\text{new}}[:,\mathrm{None}]
}.
$$

于是更新：

$$
\boxed{
\boldsymbol{l}_{\text{new}}
=\boldsymbol{\alpha}\odot\boldsymbol{l}_{\text{old}}
+\operatorname{rowsum}(\boldsymbol{P}_j^{*})
}
$$

和

$$
\boxed{
\widetilde{\boldsymbol{O}}_{\text{new}}
=\boldsymbol{\alpha}[:,\mathrm{None}]\odot
\widetilde{\boldsymbol{O}}_{\text{old}}
+\boldsymbol{P}_j^{*}\boldsymbol{V}_j
}.
$$

所有 KV tiles 处理完成后：

$$
\boxed{
\boldsymbol{O}_i
=\widetilde{\boldsymbol{O}}/\boldsymbol{l}[:,\mathrm{None}]
}.
$$

\begin{criterion}{为什么更新以后仍然精确}
新最大值只改变指数的公共参考点。把旧 $l$ 与旧 $\widetilde O$ 同时乘以 $\mathrm{e}^{m_{\text{old}}-m_{\text{new}}}$，等价于把所有旧项重新表达在新的指数基准下；随后再加入新 tile 的项。因此每一步都保持上面的三条状态定义，最终结果与整行一次性 stable softmax 完全相同。
\end{criterion}

这也是理解 FlashAttention 时最值得记住的三元组：

$$
\boxed{(\boldsymbol{m},\boldsymbol{l},\widetilde{\boldsymbol{O}})}.
$$

回到前面的母例，把前两个 score $(1,0)$ 作为第一个 tile，把 score $2$ 作为第二个 tile。处理完第一个 tile：

$$
m_{\mathrm{old}}=1,\qquad
l_{\mathrm{old}}=1+\mathrm{e}^{-1},\qquad
\widetilde{\boldsymbol{o}}_{\mathrm{old}}=(1,\mathrm{e}^{-1}).
$$

第二个 tile 到来后，新最大值变成 $2$，所以

$$
\alpha=\mathrm{e}^{1-2}=\mathrm{e}^{-1},
$$

而新 tile 的未归一化权重就是 $1$。于是

$$
l_{\mathrm{new}}
=\mathrm{e}^{-1}(1+\mathrm{e}^{-1})+1
=1+\mathrm{e}^{-1}+\mathrm{e}^{-2},
$$

$$
\widetilde{\boldsymbol{o}}_{\mathrm{new}}
=\mathrm{e}^{-1}(1,\mathrm{e}^{-1})+(1,1)
=(1+\mathrm{e}^{-1},\ 1+\mathrm{e}^{-2}).
$$

最终 $\widetilde{\boldsymbol{o}}_{\mathrm{new}}/l_{\mathrm{new}}$ 与前面整行一次性 softmax 得到的输出完全相同。这个数值轨迹直接展示了为什么旧 tile 必须在最大值变化时同时重标定 $l$ 和 $\widetilde O$。

\section{Causal tiling：能在 tile 层跳过的工作，不要等到元素层再 mask}

在 causal self-attention 中，score matrix 的上三角无效。若 query/key tile 对齐，那么对某个 query tile：

\begin{itemize}
  \item 对角线左侧的完整 KV tiles 全部合法；
  \item 对角线 tile 内需要 element-wise causal mask；
  \item 对角线右侧的完整 KV tiles 全部非法，可以直接不计算。
\end{itemize}

因此高效 kernel 常把 causal 区域拆成“完整有效 tile”与“对角边界 tile”。这里体现的是一种常见优化思路：\textbf{若粗粒度结构已经能证明一整块工作无贡献，就不要先算完再逐元素归零。}

\section{GPU 实现最小模型：软件任务层级与硬件执行资源要分开}

为了理解 FlashAttention kernel，只需要掌握一套最小 CUDA execution model，而不需要把本册变成 CUDA 教材。

\begin{itemize}
  \item kernel launch 产生一个 grid；
  \item grid 由 thread blocks 组成；
  \item NVIDIA GPU 上一个 warp 通常由 32 个 threads 组成；
  \item thread block 被调度到 streaming multiprocessor（SM）执行；
  \item 一个 SM 在资源允许时可以同时驻留多个 blocks / warps；
  \item registers、shared memory/L1、L2、device global memory 构成不同作用域和代价的存储层级；在高端数据中心 GPU 上，device global memory 常由 HBM 承载。
\end{itemize}

因此必须阻断三个错误等号：

$$
\text{thread}\neq\text{physical CUDA core},
$$

$$
\text{warp}\neq\text{thread block},
$$

$$
\text{SM}\neq\text{thread block}.
$$

\begin{handbookdiagram}
\centering
\begin{tikzpicture}[
  >=Stealth,
  font=\small,
  color=themefg,
  text=themefg,
  sw/.style={draw=themecurve,fill=themecurve!10!themebg,rounded corners=4pt,minimum width=2.0cm,minimum height=0.76cm,align=center,line width=0.9pt},
  hw/.style={draw=themeamber,fill=themeamber!10!themebg,rounded corners=4pt,minimum width=2.65cm,minimum height=0.82cm,align=center,line width=0.92pt},
  tile/.style={draw=themegreen,fill=themegreen!10!themebg,rounded corners=4pt,minimum width=2.45cm,minimum height=0.82cm,align=center,line width=0.92pt},
  flow/.style={->,draw=themefg!90,line width=1pt},
  map/.style={->,draw=themegray,dashed,line width=0.95pt},
  mem/.style={<->,draw=themeamber,line width=0.95pt}
]

\begin{pgfonlayer}{background}
  \fill[fill=themecurve!4!themebg,draw=themecurve!40,dashed,rounded corners=7pt]
    (-0.2,1.25) rectangle (9.55,5.1);
  \fill[fill=themeamber!4!themebg,draw=themeamber!40,dashed,rounded corners=7pt]
    (9.95,1.25) rectangle (16.65,5.1);
  \fill[fill=themegreen!4!themebg,draw=themegreen!40,dashed,rounded corners=7pt]
    (-0.2,-3.45) rectangle (16.65,0.75);
\end{pgfonlayer}

% CUDA software hierarchy: generic structural nesting, not a concrete launch diagram.
\node[anchor=west,font=\bfseries,color=themecurve] at (0.05,4.72) {CUDA software hierarchy};
\node[sw] (grid) at (1.2,3.05) {Grid};
\node[sw] (block) at (3.75,3.05) {Thread Block / CTA};
\node[sw] (warp) at (6.35,3.05) {Warp\\\scriptsize 32 threads};
\node[sw] (thread) at (8.55,3.05) {Thread};
\draw[flow] (grid) -- (block);
\draw[flow] (block) -- (warp);
\draw[flow] (warp) -- (thread);
\node[anchor=west,font=\scriptsize,color=themegray] at (0.25,1.72)
  {Grid 含多个 blocks；每个 block 再划分为 warps。这里画的是层级，不表示只有一个 block / warp。};

% Hardware hierarchy. Mapping is scheduling, not identity.
\node[anchor=west,font=\bfseries,color=themeamber] at (10.2,4.72) {GPU hardware};
\node[hw] (sm) at (11.55,3.2) {Streaming Multiprocessor\\\scriptsize registers · shared/L1 · warp schedulers};
\node[hw] (l2) at (14.85,3.2) {Shared L2 Cache};
\node[hw] (hbm) at (14.85,1.75) {Device DRAM / HBM};
\draw[mem] (sm) -- (l2);
\draw[mem] (l2) -- (hbm);
\draw[map] (block.north) -- ++(0,0.68) -| node[pos=0.62,above=1pt,font=\scriptsize,color=themegray]{scheduled to one SM} (sm.north west);
\node[anchor=west,font=\scriptsize,color=themegray] at (10.25,1.35)
  {一个 SM 可随资源允许驻留多个 blocks；Thread 不与“某个 CUDA core”建立一一身份关系。};

% Triton abstraction: program instance and data tile are compiler-level concepts.
\node[anchor=west,font=\bfseries,color=themegreen] at (0.05,0.35) {Triton blocked-program view};
\node[tile] (tgrid) at (1.45,-1.25) {Launch Grid};
\node[tile] (pid) at (4.45,-1.25) {Program Instance\\\scriptsize $pid=\mathrm{program\_id}$};
\node[tile] (dtile) at (7.55,-1.25) {Logical Data Tile\\\scriptsize e.g. one $Q_i$ block};
\node[hw] (lower) at (10.85,-1.25) {Compiler Lowering};
\node[hw] (exec) at (14.45,-1.25) {Warps / Threads + On-chip State\\\scriptsize registers · shared memory};
\node[tile,minimum width=3.1cm] (cfg) at (10.85,-2.65) {Execution Config\\\scriptsize num\_warps · num\_stages · tile sizes};

\draw[flow] (tgrid) -- (pid);
\draw[flow] (pid) -- (dtile);
\draw[map] (dtile) -- node[above=1pt,font=\scriptsize,color=themegray]{lowered by compiler} (lower);
\draw[map] (cfg.north) -- (lower.south);
\draw[map] (lower) -- (exec);

\node[anchor=west,font=\scriptsize,color=themegray] at (0.15,-3.12)
  {Triton 的 program instance 首先描述“一个程序实例处理哪块数据”；实际用多少 warps、threads、registers/shared memory 是编译与配置后的执行映射。};

\end{tikzpicture}

\diagramcaption{CUDA 与 Triton 的抽象层映射。上半部分把 CUDA software hierarchy 与 GPU hardware hierarchy 分开：Thread Block / CTA 被调度到某个 SM，但两者不是同一对象；下半部分展示 Triton 的 program instance / data tile 经编译器与 execution config 映射到 warps、threads 和片上状态。虚线表示实现映射，不表示对象相等。}
\end{handbookdiagram}

同一 warp 内发生 data-dependent branch divergence 时，硬件需要分别执行实际被采用的路径，并对当前路径不参与的 lanes 做 inactive masking；因此 divergence 影响吞吐。现代 NVIDIA GPU 支持 Independent Thread Scheduling，所以不能把“warp 内线程永远逐指令锁步”当成程序正确性的永久假设。

\section{Triton 的 blocked-program 抽象：先描述数据块，再让编译器映射线程}

CUDA C 常从“每个 thread 负责哪个 scalar/work item”开始；Triton 更强调一个 program instance 对一个 tensor block 的运算。

典型一维结构是：

\begin{lstlisting}[language=Python]
pid = tl.program_id(0)
offsets = pid * BLOCK_SIZE + tl.arange(0, BLOCK_SIZE)
mask = offsets < n
x = tl.load(x_ptr + offsets, mask=mask)
\end{lstlisting}

这里

$$
\texttt{BLOCK\_SIZE}
$$

描述的是 program 要处理的数据范围，不意味着“一定存在同样数量的物理 threads”。执行线程数还受编译配置与 backend 映射影响。

对 FlashAttention，一个自然 program grid 是让 program instance 对应

$$
(batch,head,query\ tile),
$$

然后固定该 query tile，在 kernel 内循环读取 KV tiles。这样不同 query tiles 可以并行，而每个 program 内部又能让当前 $Q$ 与运行状态在片上复用。

\subsection{为什么不能把 \texttt{tl.load} 定义成“HBM 到 shared memory”}

\texttt{tl.load} 的语言语义是从 pointer 指定的位置加载 tensor values。编译器最终可能利用 registers、cache、shared memory、异步数据搬运等不同机制；这些属于 lowering 与目标硬件实现。

因此更安全的模型是：

\begin{processchain}
  \processstage{Triton tensor program}
  \processstage{Pointer / layout information}
  \processstage{Compiler lowering}
  \processstage{Registers / cache / shared-memory staging}
  \processstage{GPU instructions}
\end{processchain}

这条链中的箭头表示实现层映射，不表示每次 load 必然依次经过所有列出的硬件层。

\subsection{Autotune：搜索的是资源交换，不是神奇常数}

Triton 的 autotune 允许程序员给出候选配置，例如 tile size、\texttt{num\_warps}、\texttt{num\_stages}，并指定 tuning key。对新的 key 组合，运行时评测候选配置并选择更快者。

\texttt{num\_warps} 表示每个 kernel instance 协同执行所使用的 warps 数；在 NVIDIA backend 上，8 warps 对应 $8\times32=256$ threads。\texttt{num\_stages} 指编译器在适用循环上进行 software pipelining 时采用的 stages 数；具体怎样 lower 到异步搬运、寄存器或 shared-memory staging，仍取决于 compiler/backend 与目标硬件。

增加 tile、warps 或 stages 可能提高数据复用与 memory/compute overlap，但也可能增加 registers、shared memory 占用并降低 occupancy。因此 autotuning 的目标是特定 workload + hardware 上的资源平衡，而不是发现一个跨设备永远最优的常数。

\section{Backward：本册只组合 Attention 的 VJP，不重讲通用 AutoDiff}

设

$$
\boldsymbol{S}=s\boldsymbol{Q}\boldsymbol{K}^{\mathsf T},
\qquad
s:=\frac{1}{\sqrt d},
$$

$$
\boldsymbol{P}=\operatorname{softmax}(\boldsymbol{S}),
$$

$$
\boldsymbol{O}=\boldsymbol{P}\boldsymbol{V}.
$$

为避免把“微分”与“loss 对张量的梯度”混成同一记号，下文用上横线表示 reverse-mode cotangent：

$$
\bar{\boldsymbol{O}}
:=\frac{\partial L}{\partial\boldsymbol{O}}.
$$

工程代码里常把这些量命名为 \texttt{dO}、\texttt{dQ}、\texttt{dK}、\texttt{dV}；这里的数学正文统一写成 $\bar O,\bar Q,\bar K,\bar V$。Area 50 AutoDiff 已经拥有 matmul VJP 与 softmax VJP，本册只把它们按 Attention 的 computation graph 组合。

首先：

$$
\boxed{
\bar{\boldsymbol{V}}
=\boldsymbol{P}^{\mathsf T}\bar{\boldsymbol{O}}
}
$$

以及

$$
\boxed{
\bar{\boldsymbol{P}}
=\bar{\boldsymbol{O}}\boldsymbol{V}^{\mathsf T}
}.
$$

对每一 softmax row，定义

$$
D_i
:=\sum_jP_{ij}\bar P_{ij}.
$$

则 softmax VJP 为

$$
\boxed{
\bar{\boldsymbol{S}}
=\boldsymbol{P}\odot
\left(
\bar{\boldsymbol{P}}-\boldsymbol{D}[:,\mathrm{None}]
\right)
}.
$$

最后

$$
\boxed{
\bar{\boldsymbol{Q}}
=s\,\bar{\boldsymbol{S}}\boldsymbol{K}
}
$$

以及

$$
\boxed{
\bar{\boldsymbol{K}}
=s\,\bar{\boldsymbol{S}}^{\mathsf T}\boldsymbol{Q}
}.
$$

$D_i$ 还可以改写成

$$
\boxed{
D_i
=\sum_k O_{ik}\bar O_{ik}
}.
$$

因为

$$
\sum_jP_{ij}\bar P_{ij}
=\sum_jP_{ij}
\sum_k\bar O_{ik}V_{jk}
=\sum_k\bar O_{ik}O_{ik}.
$$

因此 backward 可以先做一次 row-wise preprocess：

$$
\boldsymbol{D}
=\operatorname{rowsum}
\left(
\boldsymbol{O}\odot\bar{\boldsymbol{O}}
\right).
$$

\section{为什么 forward 保存 LogSumExp，而不是保存完整概率矩阵}

Forward 已经计算过每行的 running maximum $m_i$ 与 denominator $l_i$。定义

$$
LSE_i
:=m_i+\ln l_i.
$$

因为

$$
P_{ij}
=\frac{\mathrm{e}^{S_{ij}-m_i}}{l_i}
=\mathrm{e}^{S_{ij}-m_i-\ln l_i},
$$

所以

$$
\boxed{
P_{ij}
=\mathrm{e}^{S_{ij}-LSE_i}
}.
$$

这意味着 backward 不需要保存完整 $N\times N$ 的 $P$。只要保留 backward 所需的输入/输出与每个 query row 的少量统计量，就可以在当前 tile 上重新计算 score 和 probability，立即用于 $\bar V,\bar S,\bar Q,\bar K$，然后丢弃。具体需要保存哪些张量取决于实现接口，不能把某一版 kernel 的保存集合写成普遍定义。

这是一种明确的 compute-memory trade-off：

$$
\boxed{
\text{少保存大型 activation}
+\text{backward 重算局部 tile}
}
$$

如果重算所增加的 compute 比大型中间矩阵的 HBM write/read 更便宜，这种交换既能降低峰值显存，也可能提升实际速度。

\section{Backward 的循环为什么要按目标梯度重新分组}

每个 $\bar Q_i$ 依赖所有 KV tiles；每个 $\bar K_j,\bar V_j$ 又依赖所有 Q tiles。若在双重循环中不断更新所有部分梯度，会造成重复 device-memory 写回或复杂同步。

更好的组织方式是：

\begin{enumerate}
  \item 先预计算 $D_i=\boldsymbol{o}_i^{\mathsf T}\bar{\boldsymbol{o}}_i$；
  \item 固定一个 KV tile，遍历所有 Q tiles，在片上累计该 tile 的 $\bar K,\bar V$，最后写回；
  \item 固定一个 Q tile，遍历所有 KV tiles，在片上累计该 tile 的 $\bar Q$，最后写回。
\end{enumerate}

\begin{handbookdiagram}
\centering
\begin{tikzpicture}[
  >=Stealth,
  font=\small,
  color=themefg,
  text=themefg,
  shared/.style={draw=themeamber,fill=themeamber!10!themebg,rounded corners=4pt,minimum width=3.0cm,minimum height=0.86cm,align=center,line width=0.95pt},
  kernel/.style={draw=themecurve,fill=themecurve!10!themebg,rounded corners=5pt,minimum width=4.0cm,minimum height=1.05cm,align=center,line width=1pt},
  result/.style={draw=themegreen,fill=themegreen!10!themebg,rounded corners=4pt,minimum width=2.2cm,minimum height=0.78cm,align=center,line width=0.95pt},
  flow/.style={->,draw=themefg!90,line width=1.05pt},
  support/.style={->,draw=themegray,dashed,line width=0.95pt}
]

\node[shared] (up) at (2.3,2.85) {$O,\ \bar O$};
\node[shared] (prep) at (6.2,2.85) {$D_i=\sum_k O_{ik}\bar O_{ik}$\\\scriptsize row-wise preprocess};
\node[shared] (saved) at (10.3,2.85) {$Q,K,V,\ LSE$\\\scriptsize saved forward state};

\draw[flow] (up) -- (prep);

\node[kernel] (kv) at (4.1,0.65) {固定一个 KV tile\\遍历全部 Q tiles\\\scriptsize recompute $P$，累计 $\bar K,\bar V$};
\node[kernel] (q) at (9.15,0.65) {固定一个 Q tile\\遍历全部 KV tiles\\\scriptsize recompute $P$，累计 $\bar Q$};

\draw[support] (prep.south) -- ++(0,-0.55) -| (kv.north);
\draw[support] (prep.south) -- ++(0,-0.55) -| (q.north);
\draw[support] (saved.south) -- ++(0,-0.45) -| (kv.north east);
\draw[support] (saved.south) -- ++(0,-0.45) -| (q.north east);

\node[result] (dk) at (4.1,-1.55) {$\bar K,\ \bar V$};
\node[result] (dq) at (9.15,-1.55) {$\bar Q$};
\draw[flow] (kv) -- node[right,font=\scriptsize,color=themegray]{一次写回} (dk);
\draw[flow] (q) -- node[right,font=\scriptsize,color=themegray]{一次写回} (dq);

\node[anchor=west,font=\scriptsize,color=themegray] at (0,-2.75)
  {两条 kernel 路径使用同一梯度公式，只改变累积顺序；目标是让当前梯度 tile 尽量留在片上直到完成。};

\end{tikzpicture}

\diagramcaption{FlashAttention backward 的两种累积方向。两条主路径共享前置的 $D$ 与 LSE/recomputed $P$，但分别固定 KV tile 或 Q tile，使目标梯度尽可能在片上完成累计后再写回。}
\end{handbookdiagram}

这个循环重排不是新的梯度公式；它是在保持相同数学依赖的前提下，选择更适合 GPU data locality 的执行顺序。

\section{正确性与性能边界：必须同时回答“算得对”和“为什么快”}

\subsection{数学正确性}

FlashAttention forward 的数学正确性来自 $(m,l,\widetilde O)$ 不变量；backward 的正确性来自对同一个标准 Attention 映射应用 matmul 与 softmax 的 VJP，并在需要时重算与 forward 一致的 masked probability。由于浮点运算次序不同，优化实现不要求与朴素实现逐 bit 相同，而应在明确的 dtype 与容差下验证数值一致性。

因此验证至少分两层：

\begin{itemize}
  \item \textbf{numerical correctness}：与高精度/框架 reference 的 output、$\bar Q,\bar K,\bar V$ 在明确且与 dtype 相匹配的浮点容差内一致；
  \item \textbf{structural correctness}：causal mask、scale、shape、stride、边界 tile 与 dtype 路径必须和 reference 语义一致。
\end{itemize}

\subsection{性能不变量不是“shared memory 永远更快”}

FlashAttention 的性能优势依赖硬件、shape、dtype、kernel schedule 与实现版本。可以稳定保留的判断只有：

\begin{itemize}
  \item 避免大型中间矩阵反复写入/读出 HBM 能减少 data movement；
  \item tile 提高数据复用，但 tile 过大可能造成 register/shared-memory pressure；
  \item software pipelining 可以 overlap memory 与 compute，但增加在途状态；
  \item 更高 occupancy 不等于更高性能，Tensor Core utilization、bandwidth、instruction mix 同样重要；
  \item 最优配置依赖目标 GPU 与 workload，不能从单次实验或单一设备的结果固化为永久常数。
\end{itemize}

\section{FlashAttention 的版本演进：优化对象会随硬件改变}

FlashAttention 1 把 attention 作为 I/O-aware exact algorithm 重新设计：用 tiling、online softmax 与 recomputation 减少 HBM traffic，同时保持标准 Attention 数学语义。

FlashAttention-2 进一步指出，仅减少 I/O 还不够。它减少 non-matmul FLOPs，提高 sequence 方向的并行度，并重新划分 thread block / warp 工作，以降低 shared-memory 读写和低 occupancy 带来的损失。

FlashAttention-3 面向 Hopper，进一步利用 Tensor Memory Accelerator（TMA）、Tensor Core 的异步能力和 warp specialization，让数据搬运、矩阵乘法与 softmax 更充分重叠，并加入面向 FP8 的低精度路径。

截至 2026-08-22，官方 FlashAttention 仓库还提供 FlashAttention-4：它使用 CuTeDSL 实现，并针对 Hopper 与 Blackwell GPU 优化。FA4 的具体算子覆盖、dtype、shape 与硬件支持仍在持续演进，因此不能把“FA4”理解成一个对所有设备和所有 Attention 变体都固定可用的永久接口。

这些版本变化说明：数学映射可以保持不变，但 kernel 的瓶颈会随着硬件执行资源、数据搬运能力和编译器抽象改变。可以长期保留的设计原则是：

\begin{mentalmodel}{硬件—算法共同设计}
先找当前硬件上最昂贵的资源：是 HBM traffic、shared-memory traffic、non-matmul instruction、Tensor Core feed、exp throughput 还是 synchronization；再重新安排 tile、流水线和工作划分。硬件的相对成本一变，最优 kernel 结构也可能随之改变。
\end{mentalmodel}

\section{实际工程调用：会用高性能 Attention 不等于必须手写 kernel}

如果目标只是训练或推理一个普通 Transformer，优先使用框架已经提供的 scaled dot-product attention / fused backend，让框架根据设备、dtype、shape 和支持条件选择实现。

手写 Triton Attention 更适合以下目的：

\begin{itemize}
  \item 学习 GPU kernel 与 data movement；
  \item 研究新的 attention variant / mask / bias / layout；
  \item 现有 fused backend 无法表达当前算子语义；
  \item 对目标硬件做明确、可测量的性能研究。
\end{itemize}

否则自定义 kernel 还要承担 correctness、backward、dtype、mask、stride、corner case、benchmark、未来硬件兼容等维护成本。

\section{最小调用接口：看到什么问题先回到哪一层}

当问题出现 Attention 相关信号时，不要一上来就背公式。先判断当前问题属于哪一层：

\begin{handbooktable}{L{0.22\linewidth}L{0.31\linewidth}Y}
\toprule
\textbf{问题层} & \textbf{第一观察对象} & \textbf{典型问题}\\
\midrule
语义层 & $Q,K,V$ 与允许的信息边 & self/cross/causal 到底改变什么；某个位置能读谁\\
架构层 & Attention 与 FFN/residual/norm/position 的组合 & 为什么 Attention 不等于 Transformer；位置信息从哪里进入\\
序列扩展层 & full token interaction、recurrent state、hybrid mixer & 当前是在优化标准 Softmax，还是已经改写 sequence mixer；KDA 为什么不等于 FlashAttention\\
深度扩展层 & residual accumulation 与 cross-layer routing & AttnRes 在 layer/depth 之间选表示，为什么不等于 token attention\\
数值层 & score range 与 row-wise normalization & 为什么减最大值；online softmax 为什么仍精确\\
实现层 & 中间矩阵、tile、memory traffic & 为什么 FlashAttention 快；哪些对象不再物化到 HBM\\
反向层 & computation graph 与 saved/recomputed state & $\bar Q,\bar K,\bar V$ 从哪来；为什么保存 LSE 而不保存 $P$\\
性能层 & workload + hardware resource balance & block size、warps、stages 为什么需要调优\\
\bottomrule
\end{handbooktable}

若问题的核心已经变成 generic JVP/VJP、reverse mode 或 custom autograd contract，就停止在本册继续扩展，转交 Area 50 AutoDiff。

\section{一页压缩：从内容路由重建整套 Attention}

\begin{mentalmodel}{一句话压缩}
Attention 的核心是“按内容选择信息来源”：标准 Softmax Attention 在 token sequence 上显式建立路由，Transformer 再把这种跨位置路由与 FFN、residual、normalization 和 position 组合；长序列可以选择保持映射不变的 FlashAttention，也可以用 recurrent/linear/hybrid attention 改写 sequence mixer；AttnRes 处理的是另一条 depth 轴上的跨层表示路由，而 MoE 又在 channel/parameter 轴上选择参数子网络。现代大模型可以同时优化这三条轴。
\end{mentalmodel}

重建本册只需要回答十四个问题：

\begin{enumerate}
  \item $Q,K,V$ 分别控制“提问、匹配、载荷”的哪一部分？
  \item $1/\sqrt d$ 为什么需要控制 score scale？
  \item self、cross、causal attention 改变的是哪一种信息可达关系？
  \item multi-head 与 positional information 分别补了什么能力？
  \item Attention 为什么不等于 Transformer？
  \item FlashAttention 与 recurrent/linear attention 的根本边界是什么？
  \item KDA 为什么需要有限状态的 decay、delta correction 与 KV write，而不能理解成普通滚动平均？
  \item 从 Gated DeltaNet 到 KDA / Gated DeltaNet-2，forget、erase、write 为什么逐渐被拆成更细控制量？
  \item token-sequence routing 与 AttnRes 的 depth-wise routing 为什么是两个不同轴？
  \item Delta AttnRes 为什么改为路由 layer delta，而不是累计 hidden state？这个结论的证据边界在哪里？
  \item Kimi K3 的 sequence、depth、channel 三条优化轴分别由什么机制负责，为什么不能互相解释？
  \item 朴素 Attention 为什么会制造昂贵的 $N\times N$ memory traffic？
  \item $(m,l,\widetilde O)$ 三个状态分别保存什么不变量？
  \item 为什么 LSE 足以让 backward 重算 $P$？
  \item 为什么 kernel 的最佳 tile/warp/pipeline 配置必须跟 workload 和硬件一起看？
\end{enumerate}

\section{资料与版本边界}

本册的稳定数学与架构主线依据 Transformer、Online Softmax、FlashAttention、FlashAttention-2、FlashAttention-3、Kimi Linear 与 Attention Residuals 的原始论文；Gated DeltaNet-2 与 Delta Attention Residuals 作为 2026 年后续研究，用来说明 finite-state memory control 和 depth-routing source representation 仍在快速演化，不作为已定型结论。FlashAttention-4 与 KDA 的当前实现状态分别以官方 FlashAttention / MoonshotAI 实现为准。Kimi K3 只作为 2026 年组合 KDA、AttnRes 与 Stable LatentMoE 的当前架构实例，不承担这些机制的定义。GPU 执行语义按 NVIDIA CUDA Programming Guide 核验；Triton program、autotune、\texttt{num\_warps}、\texttt{num\_stages} 与 fused-attention 行为按 2026-08-22 的 Triton 官方文档核验。具体框架 API、backend 支持、dtype/shape 覆盖、模型参数与性能数字均属于版本敏感信息，使用时应重新核对当前官方文档。

\begin{itemize}
  \item Triton Config / autotune：\href{https://triton-lang.org/main/python-api/generated/triton.Config.html}{triton.Config}，\href{https://triton-lang.org/main/python-api/generated/triton.autotune.html}{triton.autotune}；
  \item Triton fused attention tutorial：\href{https://triton-lang.org/main/getting-started/tutorials/06-fused-attention.html}{Fused Attention}；
  \item FlashAttention 官方仓库：\href{https://github.com/Dao-AILab/flash-attention}{Dao-AILab/flash-attention}；
  \item Kimi Linear：\href{https://arxiv.org/abs/2510.26692}{Kimi Linear paper}，\href{https://github.com/MoonshotAI/Kimi-Linear}{MoonshotAI/Kimi-Linear}；
  \item KDA 当前参考实现：\href{https://github.com/MoonshotAI/minitriton/blob/main/minitriton/ops/kda.py}{MoonshotAI/minitriton KDA forward}；
  \item Attention Residuals：\href{https://arxiv.org/abs/2603.15031}{Attention Residuals}；
  \item Gated DeltaNet-2：\href{https://arxiv.org/abs/2605.22791}{Gated DeltaNet-2: Decoupling Erase and Write in Linear Attention}；
  \item Delta Attention Residuals：\href{https://arxiv.org/abs/2605.18855}{Delta Attention Residuals}；
  \item Kimi K3 技术报告：\href{https://arxiv.org/abs/2607.24653}{Kimi K3: Open Frontier Intelligence}，\href{https://github.com/MoonshotAI/Kimi-K3}{MoonshotAI/Kimi-K3}。
\end{itemize}

\end{document}
